% 本文件由 03_代码/p4_tables.py 自动生成，请勿手工修改。
% 数据源：06_支撑材料/p4_results.json

\begin{table}[htbp]
  \centering
  \caption{4-3 退款主口径与不退款敏感性口径的结果对照（两种口径分别标定）}
  \label{tab:p4-refund}
  \small
  \begin{tabular}{lrrr}
    \toprule
    指标 & 主口径（退款） & 不退款口径 & 差异（不退款 $-$ 主口径） \\
    \midrule
    合计费用 & 14,153,285.53 & 14,402,403.87 & +249,118.34 \\
    其中 计划费 & 13,802,988.74 & 13,778,495.35 & -24,493.39 \\
    其中 调整费 & 197,228.01 & 502,334.29 & +305,106.28 \\
    其中 紧急费 & 153,068.79 & 121,574.23 & -31,494.56 \\
    计划量（初始） & 21,399,492.24 & 21,344,252.60 & -55,239.64 \\
    最终常规量 & 20,842,200.33 & 21,757,381.40 & +915,181.06 \\
    上调量 & 436,206.27 & 413,128.80 & -23,077.48 \\
    下调量 & 993,498.18 & 0.00 & -993,498.18 \\
    紧急购电量 & 29,324.65 & 23,843.42 & -5,481.23 \\
    弃电量 & 1,667,711.93 & 2,601,538.42 & +933,826.49 \\
    期末储电量 & 5,889.57 & 6,099.71 & +210.14 \\
    调整提交次数 & 1,002 & 1,002 & --- \\
    \bottomrule
  \end{tabular}
  \par\vspace{2pt}
  \begin{minipage}{\textwidth}\footnotesize
    注：退款口径下，被取消部分退回原价并按 $50\%$ 计收违约费用；不退款口径下，计划购电费全额计收，$50\%$ 为额外违约费用。两种口径分别标定参数并求解，因此表中差异同时包含计费规则与参数选择的影响。两种口径下 4-3 相对 4-2 的费用关系一致。
  \end{minipage}
\end{table}
